\appendix
\section{Hyperparameters and training protocol}\label{app:hyper}
\begin{table}[h]\centering\caption{Hyperparameters, fixed across studies unless noted.}
\begin{tabular}{ll}\hline
Parameter & Value \\ \hline
Max sequence length $L$ & 150 (pilot: 60; study 14: 200) \\
PepCNN kernels / channels & widths 3/5/7, 32 channels each, 2 conv layers/branch \\
PepGNN hidden width & 32, 3 graph-conv layers, chain-window adjacency \\
Optimizer & Adam, lr $10^{-3}$ \\
Batch size & 128 (study 7); 64 (published-benchmark CNN) \\
Epochs & 8 (pilot: 10) \\
Loss & weighted BCE, $\omega = n_-/n_+$ \\
Seeds & 27 (split), 7 (training) \\
Annealing steps / restarts & 250 steps, 24 seeds, $T_0=0.2$, cooling factor 0.99 \\
Bootstrap resamples $B$ & 500 (studies 1/7); 2{,}000 (study 9) \\
Permutations $B$ & 500 \\
\hline\end{tabular}\end{table}

Training used weighted binary cross-entropy with the study-specific class
weight set by the script; the benchmark studies use balanced classes. The scale-up
training subsample of 30{,}000 sequences (15{,}000 positive) was drawn
uniformly from the train split once, then fixed. All models were evaluated on
the identical held-out test split; no test-point informed any training
decision. The pilot's 5-fold cross-validation used fold assignments stratified
by label with the same fixed seed.

\section{Data provenance}\label{app:prov}
\begin{table}[h]\centering\caption{Fetch provenance (hashes truncated to 16 hex chars as recorded).}
\begin{tabular}{lll}\hline
File & Records & Source \\ \hline
\texttt{uniprot\_kw0929\_all\_10-150.fasta} & 27{,}392 & UniProt REST, reviewed, KW-0929, len 10--150 \\
\texttt{uniprot\_nonamp\_pool\_10-150.fasta} & 105{,}041 & UniProt REST, reviewed, NOT KW-0929, len 10--150 \\
\texttt{apd6\_natural\_2024.fasta} & 3{,}306 & aps.unmc.edu APD6 downloads (2024 list) \\
\texttt{uniprot\_uncharacterized\_10-60.fasta} & 852 & UniProt REST, reviewed, uncharacterized, len 10--60 \\
\hline\end{tabular}\end{table}

Exact URLs, byte counts, SHA-256 prefixes, and UTC fetch timestamps are logged
in \texttt{data/raw/large\_provenance.json} and
\texttt{data/raw/uncharacterized\_provenance.json} in the repository. The later published-benchmark downloads have a separate SHA-256 manifest
in \texttt{results/benchmark\_sources.json}.

\section{Full design and screen listings}\label{app:full}
Table~\ref{tab:designs} lists the top 10 of 24 scale-up anneals; the remaining
14 scored between 0.877 and 0.905 with novelty 0.89--0.95 and are recorded in
\texttt{results/study07\_scaleup.json}. Table~\ref{tab:screen} lists the top
15 of the 852 screened uncharacterized proteins; the full top 25, including
CNN and GNN sub-scores, sequences, and headers, is recorded in
\texttt{results/study08\_discovery.json}. The duplicate top hit
(P0ACW4/P0ACW5) arises from two accessions carrying an identical 57-residue
sequence; we report both rather than silently de-duplicating, and flag it so
downstream users do not count it as independent evidence.

\section{Additional derivations}\label{app:math}
\begin{lemma}[Union-find clustering complexity]
Construction of the shared-$k$-mer clusters of Definition~\ref{def:cluster}
over $n$ sequences of total length $M$ runs in $O(M k\, \alpha(n))$ expected
time and $O(n + U)$ space, where $U$ is the number of distinct $k$-mers.
\end{lemma}
\begin{proof}
Each sequence contributes at most $\ell - k + 1$ $k$-mers; hashing each
$k$-mer to its first-seen sequence costs $O(1)$ expected time, so building
the map costs $O(Mk)$ expected time. Each sequence then unions with the
first-seen owner of each of its $k$-mers: at most $O(Mk)$ union operations,
each $O(\alpha(n))$ amortized by path compression with union by rank.
Storage is one parent pointer per sequence and one hash entry per distinct
$k$-mer.
\end{proof}

\begin{proposition}[Permutation test validity]
If $(y_i, s_i)$ are exchangeable under $H_0$, the Monte Carlo $p$-value
$\hat p = (1 + \sum_b \mathbf{1}[T_b \ge T_{\mathrm{obs}}])/(1+B)$ satisfies
$P(\hat p \le u) \le u$ for all $u \in [0,1]$ under $H_0$.
\end{proposition}
\begin{proof}
Under exchangeability the observed statistic and the $B$ permuted statistics
are exchangeable, so the rank of $T_{\mathrm{obs}}$ among the $B+1$ values is
uniform on $\{1,\dots,B+1\}$, with ties broken at random. Hence
$\hat p$ is uniform on $\{1/(B+1), \dots, 1\}$ up to ties, and
$P(\hat p \le u) \le u$ follows.
\end{proof}

\begin{remark}[Why not asymptotic AUROC intervals]
The DeLong estimator assumes a fixed score function and i.i.d.\ pairs; our
evaluation resamples \emph{clusters}, and cluster-level dependence violates
the i.i.d.\ assumption at the sequence level. The percentile bootstrap at the
sequence level is itself optimistic for the same reason; a cluster-level
bootstrap is the fully conservative choice and is left as an extension
noted in the repository's issue tracker.
\end{remark}
